\chapter*{General introduction}
\addcontentsline{toc}{chapter}{General introduction}
\markboth{General introduction}{General introduction}

In recent years, organisations have undergone an accelerated digital transformation, leading to an exponential increase in the production of digital documents such as technical reports, internal procedures, contracts, and operational guidelines. While this information is essential for daily decision-making, it is often dispersed across multiple systems, shared drives, and communication channels, making it difficult to access efficiently.

This fragmentation of knowledge creates a significant operational challenge. Employees frequently struggle to locate relevant information, rely on informal communication channels, or recreate existing knowledge, which results in time loss, inconsistency, and reduced productivity. This issue becomes even more critical in engineering environments where accurate and timely access to information is essential.

Beyond simple document storage, organisations now require intelligent systems capable of organising knowledge, controlling access according to organisational structure, and enabling efficient retrieval of relevant information. At the same time, the increasing use of artificial intelligence introduces new opportunities for improving how users interact with internal knowledge bases, particularly through natural language interfaces.

This final-year project was conducted during an internship at \CompanyName{}. It aims to explore and address the challenges of enterprise knowledge management by designing and implementing a platform that centralizes organisational documents and improves information accessibility through intelligent search and interaction mechanisms.

The main problem addressed in this work can be stated as follows:
\emph{How can an engineering organisation structure its internal knowledge, ensure secure and role-based access to information, and enable efficient retrieval of reliable answers from its documentation?}

To address this problem, the project focuses on the following objectives:
\begin{enumerate}[leftmargin=*,itemsep=0.2em]
  \item design a secure platform supporting multiple user roles with controlled access to organisational documents;
  \item improve information retrieval from heterogeneous document sources;
  \item enable more efficient access to knowledge through intelligent assistance mechanisms;
  \item ensure system reliability through testing and validation approaches;
  \item provide a reproducible and deployable software solution.
\end{enumerate}

The scope of this project includes the design and implementation of a full-stack web application for document management and knowledge access. It covers authentication, document organisation, and intelligent retrieval functionalities, while ensuring structured access based on organisational roles and departments.

However, certain aspects remain outside the scope of this work, including mobile application development, large-scale distributed infrastructure deployment, and external enterprise authentication systems.

This report is organised as follows. Chapter~\ref{ch:state-of-the-art} presents the state of the art and positions the project within existing solutions. Chapter~\ref{ch:requirements} describes the requirements analysis and methodology. Chapter~\ref{ch:architecture} details the system design and architecture. Chapter~\ref{ch:implementation} presents the implementation process. Chapter~\ref{ch:validation} covers testing and validation. Finally, the conclusion summarises the contributions and discusses future work.